\section{玻色子的二次量子化}

Gross 第 3 章专论玻色子。与费米子平行，但允许任意占据数，对易关系代替反对易。

\subsection{对易关系与 Fock 态}

\begin{equation}
  [b_\alpha,b_\beta^\dagger]=\delta_{\alpha\beta},
  \qquad
  [b_\alpha,b_\beta]=0.
\end{equation}
\begin{equation}
  |n_1,n_2,\ldots\rangle
  =\prod_\alpha
  \frac{(b_\alpha^\dagger)^{n_\alpha}}{\sqrt{n_\alpha!}}
  |0\rangle,
  \qquad
  n_\alpha=0,1,2,\ldots.
\end{equation}
场算符 $\psi(\mathbf{r})=\sum_\alpha\phi_\alpha(\mathbf{r})b_\alpha$ 满足
$[\psi(\mathbf{r}),\psi^\dagger(\mathbf{r}')]=\delta(\mathbf{r}-\mathbf{r}')$。

\subsection{与费米子的关键差别}

\begin{itemize}
  \item 无 Pauli 限制，可发生 Bose 凝聚（宏观占据单一轨道）；
  \item Wick 定理用对易收缩，图中无费米负号；
  \item 凝聚体需把 $b_0$ 换成 $c$-数加涨落（Bogoliubov），见第 8 章。
\end{itemize}
声子、磁振子、冷原子等均用此语言；电子--声子问题中声子部分即玻色二次量子化。
